\documentclass[11pt]{article}
\usepackage[margin=1in]{geometry}
\usepackage{amsmath,amssymb}
\usepackage{booktabs}
\usepackage{hyperref}
\usepackage{xcolor}
\usepackage{siunitx}

\title{Independent Replication --- Figueiras et al.\ (2018):\\
Split-Step (BPM) Solver for the Schr\"odinger Equation}
\author{PDE-100 Replication Project}
\date{}

\begin{document}
\maketitle

\paragraph{Paper.} E. Figueiras, D. Olivieri, A. Paredes, H. Michinel,
``An open source virtual laboratory for the Schr\"odinger equation,''
\emph{European Journal of Physics} \textbf{39} (2018) 055802.
DOI: \href{https://doi.org/10.1088/1361-6404/aac999}{10.1088/1361-6404/aac999}.
Open Access (CC-BY 3.0, IOP). PDF sha256:
\texttt{034a26a1f606e6b1f5c5a0135a89d45c0ef137b5e763cabb10a190d67e933486}.

\paragraph{Set.} PDE-100 \textperiodcentered\ Family: Schr\"odinger / nonlinear
Schr\"odinger (NLSE) --- \emph{previously uncovered in the set}.

\paragraph{Replication mode.} Fully independent from-scratch reimplementation
of the split-step Fourier beam-propagation-method (BPM) scheme using only
\texttt{numpy} FFTs. The authors' supplementary code library was
\textbf{not} downloaded or consulted; the solver was written directly from
the paper's equations. Correctness was validated against \textbf{closed-form
analytic solutions first}, then the paper's headline phenomena and stated
accuracy order were reproduced.

\section{Paper Summary}
The paper presents a simple educational Python library that integrates the
dimensionless time-dependent Schr\"odinger equation
\begin{equation}
i\,\partial_t \psi = -\tfrac{1}{2}\nabla^2 \psi + V\psi, \qquad
V = V(x)\ \text{or}\ V(\psi) \label{eq:tdse}
\end{equation}
with a first-order split-step Fourier (Lie-splitting) method,
\begin{equation}
\psi(x, t+\Delta t) = \mathcal{F}^{-1}\!\left[
e^{-i\,\Delta t\,k^2/2}\,\mathcal{F}\!\left[e^{-i\,\Delta t\,V}\psi(x,t)\right]
\right]. \label{eq:splitstep}
\end{equation}
Realized as steps I--V (potential kick in position space $\to$ FFT $\to$
kinetic phase in Fourier space $\to$ IFFT $\to$ advance $t$). Demonstrations
include reflectionless P\"oschl--Teller scattering, barriers, double wells,
2D vortex/Gaussian beams (linear); bright solitons, soliton collisions,
filamentation, BEC/Gross--Pitaevskii vortex precession (nonlinear). The
paper states the scheme conserves $\int|\psi|^2$ for real $V$ and is
first-order accurate, error $O(\Delta t)$.

\section{Claims and Results}
\begin{table}[h]
\centering\small
\begin{tabular}{@{}clll@{}}
\toprule
ID & Claim & Type & Result \\
\midrule
C1 & Eq.~\eqref{eq:splitstep} integrates TDSE; norm conserved for real $V$ & method & \textbf{Reproduced} \\
C2 & $V=-s(s{+}1)/\cosh^2 x$ reflectionless for integer $s$ (Fig.~1: $s{=}10$) & physics & \textbf{Reproduced} \\
C3 & Bright soliton eq.~(5) propagates \& two-soliton collision preserved (Fig.~2) & physics & \textbf{Reproduced} \\
C4 & First-order accuracy, error $O(\Delta t)$ & accuracy & \textbf{Reproduced} \\
\bottomrule
\end{tabular}
\end{table}

\begin{table}[h]
\centering\small
\begin{tabular}{@{}lll@{}}
\toprule
Quantity & Paper & This replication \\
\midrule
Norm conservation (real $V$) & conserved & $|{\rm norm}-1| \approx 10^{-13}$ \\
Free-propagation accuracy    & (implicit exact) & $L_2$ vs analytic $\approx 10^{-14}$; $\sigma(T)$ to 12 digits \\
Reflectionless, $s=10$ (Fig.~1) & $T=1$, $R=0$ & $R\approx 2.6\times10^{-8}$, $T=1.000000$; closed form $R=0$ \\
Bright soliton shape (eq.~5) & undisturbed & peak $1.00007$, field $L_2\approx 7\times10^{-4}$ over $T=15$ \\
Two-soliton collision (Fig.~2) & unchanged & peaks $1.0\to 0.99989$; norm $=4.0000$ \\
Order of accuracy & first order, $O(\Delta t)$ & observed 1.0005 / 1.0002 / 1.0001 \\
\bottomrule
\end{tabular}
\end{table}

\section{Method (Independent Implementation)}
\textbf{Environment:} CherryRd (local, Darwin), CPU only. Python 3.14.6;
numpy 2.4.3, scipy 1.18.0, matplotlib 3.10.8.
\begin{enumerate}
\item \texttt{work/bpm.py} --- \texttt{BPM1D}/\texttt{BPM2D} classes.
Uniform periodic grid; angular wavenumbers $k=2\pi\cdot\text{fftfreq}(N,\Delta x)$;
kinetic propagator $\exp(-i\,\Delta t\,k^2/2)$. One step = potential kick
(position) $\to$ FFT $\to$ kinetic phase (Fourier) $\to$ IFFT (paper
eq.~2 / steps I--V). Nonlinear term supplied as $V(\psi)=\kappa|\psi|^2$.
\item \texttt{test1\_free\_gaussian.py} --- exact free Schr\"odinger
propagator of a minimum-uncertainty Gaussian; IC match ($8\times10^{-17}$),
$L_2$ vs analytic at $T=8$, group velocity, spreading law.
\item \texttt{test3\_reflectionless.py} --- broad Gaussian packet
on $V=-s(s{+}1)/\cosh^2 x$; reflected fraction from Fourier partition of
$\psi$ after the packet clears the well.
\item \texttt{test3d\_verify\_formula.py} --- exact closed form
$R(k,s)=\sin^2(\pi s)/(\sinh^2(\pi k)+\sin^2(\pi s))$, so $R=0$ iff
$s\in\mathbb{Z}$ (numerically $10^{-32}$).
\item \texttt{test4\_soliton.py} --- single-soliton vs exact eq.~(5); two-soliton
head-on collision.
\item \texttt{test5\_order\_selfconv.py} --- Cauchy self-convergence on the
nonlinear soliton (non-commuting operators $\Rightarrow$ genuine Lie error).
\end{enumerate}

\section{Multi-Judge Assessment}
Free Argo endpoints, temperature 0.
\begin{table}[h]
\centering\small
\begin{tabular}{@{}lll@{}}
\toprule
Judge & Verdict & Coverage \\
\midrule
argo:gpt-5.2 & REPLICATED & C1,C2,C3,C4 \\
argo:gemini-2.5-pro & REPLICATED & C1,C2,C3,C4 \\
argo:gpt-4.1 & REPLICATED & C1,C2,C3,C4 \\
\bottomrule
\end{tabular}
\end{table}
Unanimous. Raw: \texttt{evidence/evidence\_judges.json}.

\section{Genuine Critique}
This section records honest limitations of the replication effort itself and
of the paper's claims as tested here. It is not boilerplate.

\paragraph{Scope of tested claims is narrow relative to the paper's
demonstrations.} We reproduced C1--C4: solver correctness, reflectionless
integer-$s$ P\"oschl--Teller scattering (Fig.~1), single/two bright soliton
dynamics (Fig.~2), and first-order accuracy. The paper additionally
demonstrates 2D linear and 2D nonlinear (Gross--Pitaevskii) phenomena
(vortex beams, filamentation, BEC vortex precession). Our \texttt{BPM2D}
class was implemented but the 2D headline demonstrations were not run to
full paper-figure fidelity; the ``REPLICATED'' verdict therefore applies
to the paper's core numerical claims and its 1D linear and 1D NLSE headline
figures, not to every figure in the paper.

\paragraph{Honest negative --- discarded auxiliary integrator.} A
stationary-Schr\"odinger ODE scattering integrator (\texttt{test3c}) was
numerically unstable for \emph{odd} integer $s$ (spurious $R\approx 0.44$)
because the reflectionless case requires exact cancellation of an
exponentially growing mode. It was discarded and replaced by the exact
closed form (Test~3d). This does not affect any paper claim --- both the
split-step wavepacket (Test~3) and the closed form confirm reflectionless
behavior for integer $s$ --- but it is a real replication-side failure
we chose not to hide.

\paragraph{Convention pinning was necessary.} The paper writes the kinetic
factor as $\exp(i\,\Delta t\,(2\pi k)^2/2)$ under an FFT convention where
the Laplacian eigenvalue is $-(2\pi k)^2$; with an angular-wavenumber grid
this is $\exp(-i\,\Delta t\,k^2/2)$. The sign/convention choice is not
inconsistent in the paper but had to be pinned by requiring the free case
to reproduce the exact analytic propagator (it does, to $10^{-14}$). A
reader working purely from the paper's formula literals could easily
misapply the sign; this is a minor pedagogical gap.

\paragraph{Accuracy claim tested by self-convergence, not manufactured
solution.} The observed first-order rate ($1.0005$, $1.0002$, $1.0001$)
comes from Cauchy self-convergence on the nonlinear soliton, where the
operators do not commute so genuine Lie error dominates. A stronger test
would compare to a Strang-splitting (second-order) or high-order reference
solution; we consider the self-convergence sufficient because the paper's
claim is only first order, but a fully manufactured-solution study was not
done.

\paragraph{Reflection metric is Fourier-partition, not a true asymptotic
$R$.} Test~3 measures the reflected fraction as
$\int_{k<0}|\hat\psi|^2 / \int|\hat\psi|^2$ after the packet clears the well.
This is a good proxy for a broad Gaussian, but a rigorously asymptotic
scattering coefficient would require projecting on outgoing plane-wave
states in position-space windows far from the well. The closed-form
verification (Test~3d) is what carries the reflectionless claim
quantitatively; Test~3 corroborates it.

\paragraph{Judge diversity is limited.} The three LLM judges are all Argo
endpoints and share training-data lineage; unanimous ``REPLICATED'' from
three related frontier models is weaker evidence than a disagreement-resolved
verdict across truly independent judge families. This is a project-wide
methodology limit, not a paper-specific one.

\paragraph{No performance comparison.} The paper positions itself as an
open-source educational library; we did not benchmark our reimplementation
against the paper's shipped library on runtime, memory, or usability. Our
claim is only that the numerical method behaves as advertised.

\section*{Verdict}
\textbf{REPLICATED.} The split-step Fourier (BPM) method for the linear and
nonlinear TDSE was independently reimplemented from scratch and validated
against exact analytic solutions (free Gaussian propagation to
$\sim 10^{-14}$, norm conservation to $\sim 10^{-13}$). All four testable
claims were reproduced quantitatively; three independent free-endpoint LLM
judges unanimously scored REPLICATED with full C1--C4 coverage.

\end{document}
